\section{Two compatible log minorants}\label{simp:minorants}

Every row in the repair class will satisfy $F/u\le.12$,
$\alpha/u\le1.05068$. We prove a slightly wider interface,
\begin{equation}\label{simp:wide-domain}
                       F/u\le1/5,\qquad\alpha/u\le6/5,
\end{equation}
which also absorbs finite evaluation errors. Throughout,
$g_j(S)=f(S+h_j)-p_j$ and $d_j=\mathbb Eg_j(R)$.
A positive combination of the minorants below stores only one submodular
function, even when it refers to two different heads.

\subsection{The moment inequality}

If $Y=g(R)$ for normalized monotone submodular $g$ with marginals at most
$a$, and $d=\mathbb EY$, then for every increasing $h$,
\begin{equation}\label{eq:new-active-factorial-base}
 \mathbb E[Yh(Y)]\le d\,\mathbb E[h(Y+a)],\qquad
 \mathbb E\prod_{r=0}^{s-1}(Y-ra)^+\le d^s\quad(s\ge1).
\end{equation}
To see this, fix an item order and telescope
$Y=\sum_j B_j\Delta_j(R_{<j})$. Each $\Delta_j$ is nonnegative and
decreasing. Conditional on $B_j=1$, $Y\le g(R_{-j})+a$.
A decreasing and an increasing function on a product space have
nonpositive covariance: induct on coordinates using the two-point
covariance $p(1-p)(u(1)-u(0))(v(1)-v(0))\le0$.
Consequently
\[
 \mathbb E[Yh(Y)]\le\sum_jp_j\mathbb E\Delta_j\,
                \mathbb Eh(g(R_{-j})+a)
 \le d\,\mathbb Eh(Y+a).
\]
Apply this to $h(y)=\prod_{r=1}^{s-1}(y-ra)^+$ and induct.
In particular $\mathbb EY^2\le d^2+ad$.

We shall use one fourth-moment truncation. For $x\ge7/2$,
\begin{equation}\label{eq:new-active-quartic-tail}
 x-7/2\le x(x-1)(x-2)(x-3)/47.
\end{equation}
With $v=x-4\ge-1/2$ the difference polynomial is
$v^4+10v^3+35v^2+3v+1/2$. It is positive for $v\ge0$; on
$[-1/2,0]$ it is at least $30v^2+3v+1/2\ge17/40$.
Thus for marginal bound $a\le Ap$ and $m=d/p$,
\begin{equation}\label{simp:truncated-moment}
 \mathbb E\min(Y/p,7A/2)\ge m-m^4/(47A^3).
\end{equation}

\subsection{Only two ordinary branches can be active on the repair rows}

By the deterministic private law and weighted AM--GM,
$T_0\le F+\alpha\sum_c\omega_c\theta_c<F+.37\alpha$.
On \eqref{simp:wide-domain} this is less than $.644u$, so private
mass is zero. The clipped HOME score is its singleton whenever supported,
since $F/\beta<2F<u$.
For a genuinely nonconstant supported affine
$T\le(p+kF)/R$, $T>p$ and $T\ge u$ imply $t>(R-1)/k$.
The two mixed-raw affine thresholds are $.2121441$ and $.23726325$,
both above $1/5$. The joint bound $T\le d_sp+b_sF$ similarly forces
$t>(1-d_s)/b_s>.28>1/5$. These programs cannot be nonconstant and supported.

The three remaining contraction slopes are at most $1/2$, whence
$p/u\ge.9$, $F/p\le2/9$, and $\alpha/p\le4/3<4$.
Clearing denominators proves
$(1+.5m)/1.02631852<(1+.4m)/1.00032357$ on $0\le m\le2/9$:
the difference is linear and has the required sign at its larger endpoint.
Therefore every nonconstant supported ordinary label is one of
\begin{equation}\label{simp:weak-affines}
 T_1=(p+.395494177d)/1.00024010,
 \qquad T_2=(p+.4d)/1.00032357.
\end{equation}
Their slopes are at most $2/5$, which improves the domain to
$p/u\ge23/25$, $d/p\le F/p\le5/23$, and $\alpha/p\le30/23$.
Singleton ties can use the constant singleton certificate.

\subsection{Ordinary minorants at low and high protection scales}

Write $\widehat K$ for \eqref{simp:strong-kernel} and put
$\kappa=.394805930160$, $e_0=.000238462746442$, $A_0=38/35$.
When $\alpha/p\le A_0$, use
\begin{equation}\label{eq:new-active-physical-minorant}
 \log\widehat K_{p,\alpha}(X,p+Y)
 \ge\log p+\kappa\min(Y/p,19/5)-e_0.
\end{equation}
Indeed decrease $X$ to $Y$ and protection to zero in $\mathcal H$.
For $K_0(y)=\mathcal H_{1,0}(y,1+y)$, partition at the rational
factor ties: capped breakpoints $\theta,\theta/(1-\theta)$ and
attenuated breakpoints $a^{-1}-1,a/(1-a)$, omitting infinite values.
On each interval $\log K_0(y)-\kappa y+e_0$ is concave.
Its 12 relevant endpoint comparisons up to $19/5$ have directed
margin greater than $1.2009\cdot10^{-15}$. Monotonicity handles the tail.
The envelope lemma then proves \eqref{eq:new-active-physical-minorant}.

By \eqref{simp:truncated-moment}, its expected excess over either score
in \eqref{simp:weak-affines} is at least
\[
 f(m)=\kappa[m-m^4/(47A_0^3)]-e_0-\log(1+km)+\log R.
\]
On $[0,15/158]$, an exact lower bound $f''>7/50$, rational bisection of
the sole possible derivative zero, and a directed value enclosure there
prove margins respectively greater than $9.2490055\cdot10^{-8}$ and
$2.8398941\cdot10^{-9}$. On the full $[0,5/23]$, rational arithmetic gives
\[
 f''\ge\frac{k^2}{(1+5k/23)^2}
       -\frac{12\kappa(5/23)^2}{47A_0^3}>1/10,
 \qquad f'(15/158)>0.
\]
The positive margin therefore extends to the whole needed domain.

For $\alpha/p\ge A_0$, use the simpler high-scale minorant
\begin{equation}\label{eq:wide-old-high-scale-minorant}
 \log\widehat K_{p,\alpha}(X,p+Y)
 \ge\log p+1/100+(7/20)\min(Y/p,14/3).
\end{equation}
Convexity of $\chi$ and monotonicity of $\mathcal F(r)/r$ imply that
$\widehat K$ increases with $X,V$, and with $\alpha$ when $\alpha\ge p$.
Set $p=1$, $X=Y=y$, and $\alpha=A_0$ to obtain
\[
 \log K_0(y)=\begin{cases}
 \log(1+y)+\chi(\log(A_0/(1+y))),&1+y\le A_0,\\
 \log(1+y)+\chi(\log(A_0/y))-\chi(\log((1+y)/y)),&1+y\ge A_0.
 \end{cases}
\]
This is nondecreasing. Partition $[0,14/3]$ into 467 rational intervals
of width at most $.01$. For each $[a,b]$, directed logarithms of exact
rational factor maxima give
$\log K_0(a)-1/100-(7/20)b>1/400$.
These are whole-interval inequalities by monotonicity; the last endpoint
also covers the entire tail. This proves \eqref{eq:wide-old-high-scale-minorant}.
Using \eqref{simp:truncated-moment} with $A=4/3$, the expected excess
is at least
\[
 q(m)=1/100+(7/20)[m-m^4/(3008/27)]-\log(1+(2/5)m).
\]
It decreases on $[0,5/23]$ because $(7/20)(1+2/23)<2/5$.
The upper bound $\log(1+x)\le x-x^2/2+x^3/3$ at $x=2/23$ gives
$q(5/23)\ge678030109/252528518400>1/500$.
Thus both ordinary minorants reach their selected scores.

\subsection{The NEW minorant and the common stored function}

For all $y\ge0$,
\[
 \log\min(1+y,2)\ge(9/10)y-(3/10)y^2.
\]
On $[0,1]$ the derivative difference has positive numerator
$1/10-3y/10+3y^2/5$; beyond one the quadratic is at most
$27/40<\log2$. The second-moment bound gives
\[
 \mathbb E\log\min(1+g_j(R)/p_j,2)
 \ge\frac{d_j}{p_j}\left[\frac9{10}-\frac3{10}\frac{\alpha+d_j}{p_j}\right]
 \ge\frac25\frac{d_j}{p_j}
 \ge\log\left(1+\frac25\frac{d_j}{p_j}\right)
\]
whenever $\alpha+F\le5p_j/3$. The left side is the mean of a
nonnegative monotone submodular function. Together with the constant
$\log p_j$, it is a pointwise minorant for the NEW payoff
$\log\min(f(S+h_j),2p_j)$.

Every supported ordinary label in the repair rectangle has one of the
preceding ordinary minorants or is constant. Every retained NEW candidate
has the capped-log minorant. Add their variable parts with the selected
row weights and retain their constant offsets. Positive sums, positive
scalings, and truncation preserve monotone submodularity; the resulting
single function has mean at least the weighted logarithmic score and is
pointwise at most the weighted payoff. Its variable range is at most two,
since the respective summand ranges are $\kappa(19/5)$, $(7/20)(14/3)$,
and $\log2$, and row weights sum to one. This establishes the local
interface required by the one-flow repair theorem.
